%  LaTeX support: latex@mdpi.com 
%  For support, please attach all files needed for compiling as well as the log file, and specify your operating system, LaTeX version, and LaTeX editor.

%=================================================================
\documentclass[ijgi,article,submit,pdftex,moreauthors]{Definitions/mdpi} 
% For posting an early version of this manuscript as a preprint, you may use "preprints" as the journal and change "submit" to "accept". The document class line would be, e.g., \documentclass[preprints,article,accept,moreauthors,pdftex]{mdpi}. This is especially recommended for submission to arXiv, where line numbers should be removed before posting. For preprints.org, the editorial staff will make this change immediately prior to posting.

%--------------------
% Class Options:
%--------------------

%=================================================================
% MDPI internal commands
\firstpage{1} 
\makeatletter 
\setcounter{page}{\@firstpage} 
\makeatother
\pubvolume{1}
\issuenum{1}
\articlenumber{0}
\pubyear{2022}
\copyrightyear{2022}
%\externaleditor{Academic Editor: Firstname Lastname}
\datereceived{} 
%\daterevised{} % Only for the journal Acoustics
\dateaccepted{} 
\datepublished{} 
%\datecorrected{} % Corrected papers include a "Corrected: XXX" date in the original paper.
%\dateretracted{} % Corrected papers include a "Retracted: XXX" date in the original paper.
\hreflink{https://doi.org/} % If needed use \linebreak
%\doinum{}
%------------------------------------------------------------------
% The following line should be uncommented if the LaTeX file is uploaded to arXiv.org
%\pdfoutput=1

%=================================================================
% Add packages and commands here. The following packages are loaded in our class file: fontenc, inputenc, calc, indentfirst, fancyhdr, graphicx, epstopdf, lastpage, ifthen, lineno, float, amsmath, setspace, enumitem, mathpazo, booktabs, titlesec, etoolbox, tabto, xcolor, soul, multirow, microtype, tikz, totcount, changepage, attrib, upgreek, cleveref, amsthm, hyphenat, natbib, hyperref, footmisc, url, geometry, newfloat, caption

%=================================================================
% Full title of the paper (Capitalized)
\Title{Droughts in Kalahari Desert, Makgadikgadi Pan and Okavango Delta, Botswana analyzed by mapping dryland landscapes using GMT}

% MDPI internal command: Title for citation in the left column
\TitleCitation{Droughts in Kalahari Desert, Makgadikgadi Pan and Okavango Delta, Botswana analyzed by mapping dryland landscapes using GMT}

% Author Orchid ID: enter ID or remove command
\newcommand{\orcidauthorA}{0000-0002-5759-1089} % Add \orcidA{} behind the author's name
%\newcommand{\orcidauthorB}{0000-0000-0000-000X} % Add \orcidB{} behind the author's name

% Authors, for the paper (add full first names)
\Author{Polina Lemenkova $^{1}$\orcidA{}}

%\longauthorlist{yes}

% MDPI internal command: Authors, for metadata in PDF
\AuthorNames{Polina Lemenkova}

% MDPI internal command: Authors, for citation in the left column
\AuthorCitation{Lemenkova, P.}
% If this is a Chicago style journal: Lastname, Firstname, Firstname Lastname, and Firstname Lastname.

% Affiliations / Addresses (Add [1] after \address if there is only one affiliation.)
\address{%
$^{1}$ \quad Université Libre de Bruxelles, École polytechnique de Bruxelles (EPB, Brussels Faculty of Engineering), Laboratory of Image Synthesis and Analysis, Building L, Campus de Solbosch CP131/3, Avenue Franklin D. Roosevelt 50, B-1050, Brussels, Belgium; polina.lemenkova@ulb.be or pauline.lemenkova@gmail.com}

% Contact information of the corresponding author
\corres{Correspondence: polina.lemenkova@ulb.be; Tel.: +32471860459 (P.L.)}

% Abstract (Do not insert blank lines, i.e. \\) 
\abstract{Botswana has arid climate, prone to droughts. It topography includes unique areas of Kalahari Desert, Makgadikgadi Pan and Okavango Delta with varying climate and environmental settings. This paper presents mapping of climate variability in Botswana. Data were taken from TerraClimate. Variables include T°Cmin, T°Cmax, precipitation, soil moisture, evapotranspiration (PET and AET), downward surface shortwave radiation, vapor pressure and vapor pressure deficit (VPD), wind speed and Palmer Drought Severity Index (PDSI). The paper describes a scripting method of automated console-based scripting using GMT by machine learning techniques in a cartographic workflow. Several GMT modules were used to derive a set of climate parameters for Botswana. The data have been supplemented with adjusted cartographic elements and inspected by GDAL. The PDSI in Botswana in 2018 shows stepwise variation with seven areas of drought: 1) -3.7 to -2.2. (extreme); 2) -2.2 to -0.8 (strong, southern Kalahari); 3) -0.8 to 0.7 (significant, central Kalahari; 4) 0.7 to 2.1 (moderate); 5) 2.1 to 3.5 (lesser); 6) 3.5 to 4.9 (low); 7) 4.9 to 6.4 (least). VPD has a general trend towards SW (Kalahari Desert, up to 3.3) and lower in the NE of Botswana (up to 1.4). Other values vary respectively for each map. The results include 12 new climate and environmental inventory maps on Botswana for 2018. }

% Keywords
\keyword{Africa; Botswana; cartography; climate; computer science; data science; drought; GMT; mapping; programming} 

% The fields PACS, MSC, and JEL may be left empty or commented out if not applicable
%\PACS{J0101}
%\MSC{}
%\JEL{}

%%%%%%%%%%%%%%%%%%%%%%%%%%%%%%%%%%%%%%%%%%
\begin{document}

%%%%%%%%%%%%%%%%%%%%%%%%%%%%%%%%%%%%%%%%%%
\section{Introduction}

Botswana is a topographically flat country with unique extra arid climate caused by the presence of Kalahari Desert which occupies up to 70\% of the territory of country (Fig. 1). The biochemical composition of the sandstones from Kalahari Desert shows a distribution of quartzose sediments and sandstones prone to chemical weathering producing high volumes of clay minerals (Franchi et al. 2021). Together with extreme degrees of temperature, such geologic and environmental setting of Kalahari limit the distribution of vegetation and controls the types of dominating plants and species adopted to such conditions. Another notable landform of Botswana is the Okavango Delta located in the NW of the country (Fig. 1), which is one of the world's largest inland deltas losing about 98\% of its water through evapotranspiration (Moses, Hambira, 2018). The Makgadikgadi Pan salt pan with specific dune, barchan and crescentic sand landforms, located over the dry savanna of NE of the country, is one of the largest salt flats in the world (Burrough et al. 2012).

Due to the extreme environmental and climate setting such as drought (Hemming, 1980; Karikari, Tabona, 2004), Botswana is one of the world’s least populated countries. Botswana people need to adapt to special living setting and adjust their social activities constrained by the environmental factors: livestock diseases, human-wildlife conflicts, drought, water supply deficit, resource scarcity, fragmented landscapes, bush encroachment of savannah ecosystems (Badimo et al. 2020; Basupi et al. 2019; Charis et al. 2019). Therefore, special successful measures and actions were introduced as state subsidies for drought relief management aimed at food and nutrition security (Belbase, Morgan, 1994). The assessment of social vulnerability to natural hazards in Botswana acts as a measure for reduction of climate and environmental disaster risks (Dintwa et al. 2019). Byakatonda et al. (2018a) recognized that hydrological droughts in Botswana may have several months of lag after the climatological droughts which shows the complex non-linear correlation between the hydrological and environmental factors associated with climate variability. 

Interest in climate assessment, the occurrence and distribution of droughts and climatic extremes in Botswana has been driven by the need to analyze environmental issues of the country: range degradation (Vanderpost et al., 2011), evolution and adaptation of agro-pastoral livelihoods to recurrent droughts (Mogotsi et al. 2013), land change in the dryland agricultural landscapes of Botswana involving expanding bare lands and rocky surfaces (Akinyemi, Mashame, 2018), biochemical analysis, e.g. studies of the effect of evapotranspiration on the evolution of dissolved inorganic carbon in the NW Botswana (Atekwana et al. 2016). 

A variety of environmental studies in Botswana focused on the analysis of rainfall repeatability, intensity and consequences in different regions of the country. Modeling the meteorological data enables to perform both retrospective and predictive assessment of the climate change variations. For instance, Kenabatho et al. (2012) performed retrospective analysis of climate effects in arid environment of Botswana to determine changes in rainfall series and identify correlations between rainfall and temperature. In contrast, a prognosis study estimating the beginning and end of summer rainfalls is presented by Byakatonda et al. (2019). Rainfall variability and trends for analysis of their effects on vulnerable dryland agriculture through climate variation were presented by Batisani, Yarnal (2010).

To contribute to the existing studies of climate in Botswana the existing research should be continued through the updated maps based on the actual datasets. In view of this the aim of this paper is to produce a series of the new climatic inventory maps of Botswana. The specific goal of this manuscript is to provide the environmental and climate context in Botswana through the advanced cartographic visualization. It involves the analysis regarding the distribution of temperature extremes in Botswana, occurrence of drought events, variation in the potential and actual evapotranspiration visualized on the maps. 

The analyzed correlations between the soil moisture and precipitation, climate variables and topographic significant landforms can be used for land management, crop monitoring and nature conservation in Botswana. The main objectives of the presented research are multifold:

    i. to produce a thematic series of climate maps of Botswana and their associated topographic, environmental and climate attributes, referenced to the regional and local geospatial geographic scale of Botswana; 
    ii. to perform scripting based mapping of the results to present a new methodology of the GMT based mapping using coding approach which in contrast to the traditional GIS applies the synthesis of the cartographic methods and programming syntax;
    iii. to apply the open source datasets of GEBCO and TerraClimate raster data with a spatial extent of Botswana to address and visualize key questions regarding the occurrence, trends and strength of different climate variables associated with environmental versus topographic factors;
    iv. to visualize environmental and climate parameters of Botswana for the year 2018 using GMT scripting techniques: temperature extremes Tmin/Tmax, precipitation, soil moisture, potential evapotranspiration (PET), actual evapotranspiration (AET), downward surface shortwave radiation, vapor pressure (VP), vapor pressure deficit (VPD), wind speed, Palmer Drought Severity Index (PDSI).
    
    Automated methods and independent modules of GMT used to process climate datasets for mapping climate and environmental variables of Botswana included reference to the cartographic workflow of GMT described in the existing technical literature (Lemenkova, 2019a, 2019b, 2021a) including references cited in the following chapters. 

%%%%%%%%%%%%%%%%%%%%%%%%%%%%%%%%%%%%%%%%%%
\section{Materials and Methods}

The methodology of recent environmental studies is mostly focused on utilizing the traditional GIS approaches associated with the menu-based Graphical User Interface (GUI) software (e.g. Kolawole et al. 2014; Klaučo et al. 2014, 2017; Schenke, Lemenkova, 2008; Suetova et al. 2005a, 2005b; Lemenkova et al. 2012), while this study is fully based on the Generic Mapping Tools (GMT) toolset (Wessel et al. 2019) which performs mapping through running the scripts via the command line from the console. This yields print-quality maps based on the raster and vector layers processed using a combination of various GMT modules as required in scripting. GMT has gained in popularity in topographic, marine geologic and geophysical mapping where highly detailed and accurate data for modelling the phenomena of the Earth processes and structures are essential (Gauger et al. 2007; Lemenkova, 2020a; Gohl et al. 2006a, 2006b; Kuhn et al. 2006).  

\subsection{Data}
	A data set of GEBCO (Schenke, 2016) was used, which includes unprecedentedly high resolution (15 arc-second) topographic and bathymetric data covering the Earth's relief surface: https://www.gebco.net/. Topographic data quality is inherently influenced by the algorithms of data capture, processing and modeling as shown earlier (Smith, 1993; Lemenkova, 2021b, 2020d; Bracke et al. 2019). The GEBCO data was thus well suited to reliably visualize the topography and geomorphic structures in the processed maps (Lemenkova, 2020b, 2020c). In addition, a dataset of TerraClimate of Botswana has been applied (Abatzoglou et al. 2018). 
	
\subsection{Mapping Methodology}
	The ability of GMT to visualize either the cartographic elements or apply the color palettes to the raster image surfaces is dependent on the adjusted settings (‘flags’) of that particular module (Lemenkova, 2020e). 
	Provided each module of GMT results in added cartographic element (e.g. ‘grdimage’ for visualizing raster image, ‘grdcontour’ for adding isolines, ‘psscale’ for adding color legend, ‘psbasemap’ for adding scale bar, directional rose, etc), GMT is capable of visualizing raster and vector data with each parameter adjusted in a line of the code written using GMT syntax. To be more illustrative, examples of the code snippets are provided below with explanations.
	Fig. 1 shows the topographic map of Botswana. The visualization has been made using a sequence of the GMT codes of which the most important are the following ones:

\begin{enumerate}
	\item	Making background raster grid with 50\% translucency: $‘gmt grdimage bw_relief.nc -Cpauline.cpt -R19.5/29.5/-27/-17.5 -JM6.5i -I+a15+ne0.75 -t50 -Xc -P -K > \$ps’$
	\item Adding isolines: $‘gmt grdcontour bw_relief1.nc -R -J -C200 -A200+f7p,26,\\darkbrown -Wthinner,darkbrown -O -K >> \$ps’$
	\item Adding coastlines, country borders and rivers: $‘gmt pscoast -R -J -Ia/thinner,blue -Na -N1/thicker,red -W0.1p -Df -O -K >> \$ps’$
	\item Adding text for cities: $‘gmt pstext -R -J -N -O -K -F+f13p,0,black+jLB -Gwhite@60 >> \$ps << EOF 25.0 -24.57 Gaborone EOF’$, where the plotted text is located between the two EOF (which is a common abbreviation for the End Of File expression).
	\item Changing text parameters: $'gmt pstext -R -J -N -O -K -F+f10p,23,blue2+jLB >> \$ps << EOF 24.5 -20.45 Makgadikgadi 24.6 -20.62 (salt pans) EOF'.$
	\item Adding a GMT logo: $‘gmt logo -Dx7.3/-2.8+o0.1i/0.1i+w2c -O -K >> \$ps’$
	\item Adding a subtitle: $‘gmt pstext -R0/10/0/15 -JX10/10 -X0.5c -Y11.8c -N -O -F+f10p,0,black+jLB >> \$ps << EOF 2.0 9.0 Digital elevation data: \\GEBCO/SRTM, 15 arc sec (ca. 450 m) resolution grid EOF’.$
	\item Plotting the insert map (round globe in the lower right bottom corner of the map in Fig. 1 showing the country’s location with the map of Africa): $‘gmt pscoast --MAP\_GRID\_PEN\_PRIMARY=thinnest,white -Rg -JG28.0/-2.0S/\$w -Da -Glightgoldenrod1 -A5000 -Bga -Wfaint -EBW+gred -Sdodgerblue -O -K -X\$x0 -Y\$y0 >> \$ps’.$ This code utilizes the ‘BW’ abbreviation of Botswana from the ISO 3166-1 alpha-2 countries codes.
	\item Specifying the location of the insert map on the main map in Fig. 1: $‘gmt psbasemap -R -J -O -K -DjBR+w3.4c+o-0.2c/-0.2c+stmp >> \$ps’$. Here the ‘BR’ signifies the ‘bottom right’ specification for the insert globe map plotting.
	\item $‘read x0 y0 w h < tmp’$ here the temporary file with the data is created.
	\item $‘gmt psxy -R -J -O -K -T  -X-\${x0} -Y-\${y0} >> \$ps’$: here the plotting is finished, this part of the script is ready and ‘-O -K’ signifies the overlay and code continuation. 
\end{enumerate}

Using the same approach of scripting presented above, all the other climate maps have been plotted suing the datasets from the TerraClimate and visualized in the Mercator projection for all the maps for compatibility reasons. The important code snippets for climate maps: 

\begin{itemize}
	\item	Extracting a subset of AET for Botswana: $‘gmt grdcut TerraClimate\_aet_2018.nc \\-R19.5/29.5/-27/-17.5 -Gbw\_aet.nc’.$
	\item Adding grids for all the maps: $‘-Bpxg1f0.5a1 -Bpyg1f0.5a1 -Bsxg1 -Bsyg1’.$ Here the ‘f’ stands for the frequency of minor ticks, ‘g’ for major and ‘a’ for annotations.
	\item Converting PostScript file to the image JPG file with 720 dpi resolution using GhostScript (PET of Botswana): $‘gmt psconvert BW_PET.ps -A0.5c -E720 -Tj -Z’.$ 
	\item Generating color palette and adjusting it for the extent of values in PDSI map of Botswana: $‘gmt makecpt -Cqual-mixed-12.cpt -V -T-5.1/12.1 > pauline.cpt’$.
	\item Using the GDAL library for analysis of the statistics for Tmin raster grid of Botswana: $‘gdalinfo -stats bw\_tim.nc’.$ 
\end{itemize}
  
    The general climate maps (Tmin/max, soil moisture, precipitation, downward surface shortwave radiation, vapor pressure, vapor pressure deficit and wind speed) and the derived reference indices (PDSI, PET, AET) used for Botswana originated from TerraClimate open source dataset for the year 2018. All the maps were plotted in the Mercator projection in order to achieve relevant matching coordinate reference grid at the identical spatial extent and congruent resolution in each map for compatibility and agreement of the map series. The maps were adjusted for the extent of Botswana and surrounding regions (neighboring countries) by subtracting the grids from the corresponding NetCDF files using ‘-R’ module, e.g.: $'gmt grdcut TerraClimate\_tmin\_2018.nc -R19.5/29.5/-27/-17.5 -Gbw\_tim.nc'.$ Here the '-R19.5/29.5/-27/-17.5' stands for the coordinates of the atusy area in WESN (West-East-South-North) convention with coordinates of southern hemisphere as negative values.
    
    This paper accepts the Palmer Drought Severity Index (PDSI) for visualizing the strength of drought over Botswana. Besides the PDSI, there are other types of indices used for meteorological and climate research, for instance: Standardized Precipitation Index (SPI), Standardized Precipitation Evapotranspiration Index (SPEI) to mention a few (Byakatonda et al. 2018b, 2018c). The PDSI was chosen since this is a reliable and reputable index known since its development by Palmer in 1965 (Palmer, 1965) based on a measurement of dryness based on recent precipitation and temperature. 

%%%%%%%%%%%%%%%%%%%%%%%%%%%%%%%%%%%%%%%%%%
\section{Results}

The text continues here. 

\subsection{Figures, Tables and Schemes}

All figures and tables should be cited in the main text as Figure~\ref{fig1}, Table~\ref{tab1}, Table~\ref{tab2}, etc.

\begin{figure}[H]
\includegraphics[width=10.5 cm]{Definitions/logo-mdpi}
\caption{This is a figure. Schemes follow the same formatting. If there are multiple panels, they should be listed as: (\textbf{a}) Description of what is contained in the first panel. (\textbf{b}) Description of what is contained in the second panel. Figures should be placed in the main text near to the first time they are cited. A caption on a single line should be centered.\label{fig1}}
\end{figure}   
\unskip

\begin{table}[H] 
\caption{This is a table caption. Tables should be placed in the main text near to the first time they are~cited.\label{tab1}}
\newcolumntype{C}{>{\centering\arraybackslash}X}
\begin{tabularx}{\textwidth}{CCC}
\toprule
\textbf{Title 1}	& \textbf{Title 2}	& \textbf{Title 3}\\
\midrule
Entry 1		& Data			& Data\\
Entry 2		& Data			& Data\\
\bottomrule
\end{tabularx}
\end{table}
\unskip

\begin{table}[H]
\caption{This is a wide table.\label{tab2}}
	\begin{adjustwidth}{-\extralength}{0cm}
		\newcolumntype{C}{>{\centering\arraybackslash}X}
		\begin{tabularx}{\fulllength}{CCCC}
			\toprule
			\textbf{Title 1}	& \textbf{Title 2}	& \textbf{Title 3}     & \textbf{Title 4}\\
			\midrule
			Entry 1		& Data			& Data			& Data\\
			Entry 2		& Data			& Data			& Data \textsuperscript{1}\\
			\bottomrule
		\end{tabularx}
	\end{adjustwidth}
	\noindent{\footnotesize{\textsuperscript{1} This is a table footnote.}}
\end{table}

%\begin{listing}[H]
%\caption{Title of the listing}
%\rule{\columnwidth}{1pt}
%\raggedright Text of the listing. In font size footnotesize, small, or normalsize. Preferred format: left aligned and single spaced. Preferred border format: top border line and bottom border line.
%\rule{\columnwidth}{1pt}
%\end{listing}

Text.

Text.

\subsection{Formatting of Mathematical Components}

This is the example 1 of equation:
\begin{linenomath}
\begin{equation}
a = 1,
\end{equation}
\end{linenomath}
the text following an equation need not be a new paragraph. Please punctuate equations as regular text.
%% If the documentclass option "submit" is chosen, please insert a blank line before and after any math environment (equation and eqnarray environments). This ensures correct linenumbering. The blank line should be removed when the documentclass option is changed to "accept" because the text following an equation should not be a new paragraph.

This is the example 2 of equation:
\begin{adjustwidth}{-\extralength}{0cm}
\begin{equation}
a = b + c + d + e + f + g + h + i + j + k + l + m + n + o + p + q + r + s + t + u + v + w + x + y + z
\end{equation}
\end{adjustwidth}

% Example of a page in landscape format (with table and table footnote).
%\startlandscape
%\begin{table}[H] %% Table in wide page
%\caption{This is a very wide table.\label{tab3}}
%	\begin{tabularx}{\textwidth}{CCCC}
%		\toprule
%		\textbf{Title 1}	& \textbf{Title 2}	& \textbf{Title 3}	& \textbf{Title 4}\\
%		\midrule
%		Entry 1		& Data			& Data			& This cell has some longer content that runs over two lines.\\
%		Entry 2		& Data			& Data			& Data\textsuperscript{1}\\
%		\bottomrule
%	\end{tabularx}
%	\begin{adjustwidth}{+\extralength}{0cm}
%		\noindent\footnotesize{\textsuperscript{1} This is a table footnote.}
%	\end{adjustwidth}
%\end{table}
%\finishlandscape

% Example of a figure that spans the whole page width. The same concept works for tables, too.
\begin{figure}[H]
\begin{adjustwidth}{-\extralength}{0cm}
\centering
\includegraphics[width=13.5cm]{Definitions/logo-mdpi}
\end{adjustwidth}
\caption{This is a wide figure.\label{fig2}}
\end{figure}  

Please punctuate equations as regular text. Theorem-type environments (including propositions, lemmas, corollaries etc.) can be formatted as follows:
%% Example of a theorem:
\begin{Theorem}
Example text of a theorem.
\end{Theorem}

The text continues here. Proofs must be formatted as follows:

%% Example of a proof:
\begin{proof}[Proof of Theorem 1]
Text of the proof. Note that the phrase ``of Theorem 1'' is optional if it is clear which theorem is being referred to.
\end{proof}
The text continues here.

%%%%%%%%%%%%%%%%%%%%%%%%%%%%%%%%%%%%%%%%%%
\section{Discussion}

Authors should discuss the results and how they can be interpreted from the perspective of previous studies and of the working hypotheses. The findings and their implications should be discussed in the broadest context possible. Future research directions may also be highlighted.

%%%%%%%%%%%%%%%%%%%%%%%%%%%%%%%%%%%%%%%%%%
\section{Conclusions}

This section is not mandatory, but can be added to the manuscript if the discussion is unusually long or complex.

%%%%%%%%%%%%%%%%%%%%%%%%%%%%%%%%%%%%%%%%%%
\section{Patents}

This section is not mandatory, but may be added if there are patents resulting from the work reported in this manuscript.

%%%%%%%%%%%%%%%%%%%%%%%%%%%%%%%%%%%%%%%%%%
\vspace{6pt} 

%%%%%%%%%%%%%%%%%%%%%%%%%%%%%%%%%%%%%%%%%%
%% optional
%\supplementary{The following supporting information can be downloaded at:  \linksupplementary{s1}, Figure S1: title; Table S1: title; Video S1: title.}

% Only for the journal Methods and Protocols:
% If you wish to submit a video article, please do so with any other supplementary material.
% \supplementary{The following supporting information can be downloaded at: \linksupplementary{s1}, Figure S1: title; Table S1: title; Video S1: title. A supporting video article is available at doi: link.}

%%%%%%%%%%%%%%%%%%%%%%%%%%%%%%%%%%%%%%%%%%
\authorcontributions{For research articles with several authors, a short paragraph specifying their individual contributions must be provided. The following statements should be used ``Conceptualization, X.X. and Y.Y.; methodology, X.X.; software, X.X.; validation, X.X., Y.Y. and Z.Z.; formal analysis, X.X.; investigation, X.X.; resources, X.X.; data curation, X.X.; writing---original draft preparation, X.X.; writing---review and editing, X.X.; visualization, X.X.; supervision, X.X.; project administration, X.X.; funding acquisition, Y.Y. All authors have read and agreed to the published version of the manuscript.'', please turn to the  \href{http://img.mdpi.org/data/contributor-role-instruction.pdf}{CRediT taxonomy} for the term explanation. Authorship must be limited to those who have contributed substantially to the work~reported.}

\funding{Please add: ``This research received no external funding'' or ``This research was funded by NAME OF FUNDER grant number XXX.'' and  and ``The APC was funded by XXX''. Check carefully that the details given are accurate and use the standard spelling of funding agency names at \url{https://search.crossref.org/funding}, any errors may affect your future funding.}

\institutionalreview{In this section, you should add the Institutional Review Board Statement and approval number, if relevant to your study. You might choose to exclude this statement if the study did not require ethical approval. Please note that the Editorial Office might ask you for further information. Please add “The study was conducted in accordance with the Declaration of Helsinki, and approved by the Institutional Review Board (or Ethics Committee) of NAME OF INSTITUTE (protocol code XXX and date of approval).” for studies involving humans. OR “The animal study protocol was approved by the Institutional Review Board (or Ethics Committee) of NAME OF INSTITUTE (protocol code XXX and date of approval).” for studies involving animals. OR “Ethical review and approval were waived for this study due to REASON (please provide a detailed justification).” OR “Not applicable” for studies not involving humans or animals.}

\informedconsent{Any research article describing a study involving humans should contain this statement. Please add ``Informed consent was obtained from all subjects involved in the study.'' OR ``Patient consent was waived due to REASON (please provide a detailed justification).'' OR ``Not applicable'' for studies not involving humans. You might also choose to exclude this statement if the study did not involve humans.

Written informed consent for publication must be obtained from participating patients who can be identified (including by the patients themselves). Please state ``Written informed consent has been obtained from the patient(s) to publish this paper'' if applicable.}

\dataavailability{In this section, please provide details regarding where data supporting reported results can be found, including links to publicly archived datasets analyzed or generated during the study. Please refer to suggested Data Availability Statements in section ``MDPI Research Data Policies'' at \url{https://www.mdpi.com/ethics}. If the study did not report any data, you might add ``Not applicable'' here.} 

\acknowledgments{In this section you can acknowledge any support given which is not covered by the author contribution or funding sections. This may include administrative and technical support, or donations in kind (e.g., materials used for experiments).}

\conflictsofinterest{Declare conflicts of interest or state ``The authors declare no conflict of interest.'' Authors must identify and declare any personal circumstances or interest that may be perceived as inappropriately influencing the representation or interpretation of reported research results. Any role of the funders in the design of the study; in the collection, analyses or interpretation of data; in the writing of the manuscript; or in the decision to publish the results must be declared in this section. If there is no role, please state ``The funders had no role in the design of the study; in the collection, analyses, or interpretation of data; in the writing of the manuscript; or in the decision to publish the~results''.} 

%%%%%%%%%%%%%%%%%%%%%%%%%%%%%%%%%%%%%%%%%%
%% Optional
\sampleavailability{Samples of the compounds ... are available from the authors.}

%% Only for journal Encyclopedia
%\entrylink{The Link to this entry published on the encyclopedia platform.}

\abbreviations{Abbreviations}{
The following abbreviations are used in this manuscript:\\

\noindent 
\begin{tabular}{@{}ll}
MDPI & Multidisciplinary Digital Publishing Institute\\
DOAJ & Directory of open access journals\\
TLA & Three letter acronym\\
LD & Linear dichroism
\end{tabular}
}

%%%%%%%%%%%%%%%%%%%%%%%%%%%%%%%%%%%%%%%%%%
%% Optional
\appendixtitles{no} % Leave argument "no" if all appendix headings stay EMPTY (then no dot is printed after "Appendix A"). If the appendix sections contain a heading then change the argument to "yes".
\appendixstart
\appendix
\section[\appendixname~\thesection]{}
\subsection[\appendixname~\thesubsection]{}
The appendix is an optional section that can contain details and data supplemental to the main text---for example, explanations of experimental details that would disrupt the flow of the main text but nonetheless remain crucial to understanding and reproducing the research shown; figures of replicates for experiments of which representative data are shown in the main text can be added here if brief, or as Supplementary Data. Mathematical proofs of results not central to the paper can be added as an appendix.

\begin{table}[H] 
\caption{This is a table caption.\label{tab5}}
\newcolumntype{C}{>{\centering\arraybackslash}X}
\begin{tabularx}{\textwidth}{CCC}
\toprule
\textbf{Title 1}	& \textbf{Title 2}	& \textbf{Title 3}\\
\midrule
Entry 1		& Data			& Data\\
Entry 2		& Data			& Data\\
\bottomrule
\end{tabularx}
\end{table}

\section[\appendixname~\thesection]{}
All appendix sections must be cited in the main text. In the appendices, Figures, Tables, etc. should be labeled, starting with ``A''---e.g., Figure A1, Figure A2, etc.

%%%%%%%%%%%%%%%%%%%%%%%%%%%%%%%%%%%%%%%%%%
\begin{adjustwidth}{-\extralength}{0cm}
%\printendnotes[custom] % Un-comment to print a list of endnotes

\reftitle{References}

% Please provide either the correct journal abbreviation (e.g. according to the “List of Title Word Abbreviations” http://www.issn.org/services/online-services/access-to-the-ltwa/) or the full name of the journal.
% Citations and References in Supplementary files are permitted provided that they also appear in the reference list here. 

%=====================================
% References, variant A: external bibliography
%=====================================
%\bibliography{your_external_BibTeX_file}

%=====================================
% References, variant B: internal bibliography
%=====================================
\begin{thebibliography}{999}
% Reference 1
\bibitem[Author1(year)]{ref-journal}
Author~1, T. The title of the cited article. {\em Journal Abbreviation} {\bf 2008}, {\em 10}, 142--149.
% Reference 2
\bibitem[Author2(year)]{ref-book1}
Author~2, L. The title of the cited contribution. In {\em The Book Title}; Editor 1, F., Editor 2, A., Eds.; Publishing House: City, Country, 2007; pp. 32--58.
% Reference 3
\bibitem[Author3(year)]{ref-book2}
Author 1, A.; Author 2, B. \textit{Book Title}, 3rd ed.; Publisher: Publisher Location, Country, 2008; pp. 154--196.
% Reference 4
\bibitem[Author4(year)]{ref-unpublish}
Author 1, A.B.; Author 2, C. Title of Unpublished Work. \textit{Abbreviated Journal Name} year, \textit{phrase indicating stage of publication (submitted; accepted; in press)}.
% Reference 5
\bibitem[Author5(year)]{ref-communication}
Author 1, A.B. (University, City, State, Country); Author 2, C. (Institute, City, State, Country). Personal communication, 2012.
% Reference 6
\bibitem[Author6(year)]{ref-proceeding}
Author 1, A.B.; Author 2, C.D.; Author 3, E.F. Title of presentation. In Proceedings of the Name of the Conference, Location of Conference, Country, Date of Conference (Day Month Year); Abstract Number (optional), Pagination (optional).
% Reference 7
\bibitem[Author7(year)]{ref-thesis}
Author 1, A.B. Title of Thesis. Level of Thesis, Degree-Granting University, Location of University, Date of Completion.
% Reference 8
\bibitem[Author8(year)]{ref-url}
Title of Site. Available online: URL (accessed on Day Month Year).
\end{thebibliography}

% If authors have biography, please use the format below
%\section*{Short Biography of Authors}
%\bio
%{\raisebox{-0.35cm}{\includegraphics[width=3.5cm,height=5.3cm,clip,keepaspectratio]{Definitions/author1.pdf}}}
%{\textbf{Firstname Lastname} Biography of first author}
%
%\bio
%{\raisebox{-0.35cm}{\includegraphics[width=3.5cm,height=5.3cm,clip,keepaspectratio]{Definitions/author2.jpg}}}
%{\textbf{Firstname Lastname} Biography of second author}

% For the MDPI journals use author-date citation, please follow the formatting guidelines on http://www.mdpi.com/authors/references
% To cite two works by the same author: \citeauthor{ref-journal-1a} (\citeyear{ref-journal-1a}, \citeyear{ref-journal-1b}). This produces: Whittaker (1967, 1975)
% To cite two works by the same author with specific pages: \citeauthor{ref-journal-3a} (\citeyear{ref-journal-3a}, p. 328; \citeyear{ref-journal-3b}, p.475). This produces: Wong (1999, p. 328; 2000, p. 475)

%%%%%%%%%%%%%%%%%%%%%%%%%%%%%%%%%%%%%%%%%%
%% for journal Sci
%\reviewreports{\\
%Reviewer 1 comments and authors’ response\\
%Reviewer 2 comments and authors’ response\\
%Reviewer 3 comments and authors’ response
%}
%%%%%%%%%%%%%%%%%%%%%%%%%%%%%%%%%%%%%%%%%%
\end{adjustwidth}
\end{document}

